% ============ 第 1 讲 课程导论 ============
\chap{第 1 讲　课程导论}{Introduction to the Course}

\textbf{本讲导读}：本讲对应 Lecture 1（Introduction），共 52 页。内容分为三块：
(1) 贸易研究的对象与问题（topics）；(2) 贸易研究的方法论（methods：模型的结构、比较静态、假设的作用、局部均衡与一般均衡）；(3) 国际贸易的基本概念与五个流行误解（misconceptions）及其现实对照。
本讲不涉及具体贸易模型，但给出的``模型方法论''是全课程每一类模型的分析框架，复习时应作为总纲掌握。

\secn{1.1 什么是国际贸易：现象与议题}

\begin{defbox}[基本定义（PPT 第 33 页，红字为原文强调）]
\begin{itemize}
  \item \term{出口（Exports）}：一国（``home''）向另一国（``foreign''）出售的产品/服务。
  \item \term{进口（Imports）}：一国从另一国购买的产品。
  \item \term{贸易顺差 / 逆差（Trade surplus / deficit）}：出口减进口的净市场价值（顺差）；反之（进口减出口）为逆差。
  \item \term{外商直接投资（FDI）}：对海外关联企业（overseas affiliates）的投资与控制。
  \item 注意：本课程\redkw{不涉及汇率}——汇率与国际收支失衡属于``国际金融''的内容。
\end{itemize}
\end{defbox}

\textbf{国际贸易的现象（第 7 页）}，四个层次：
\begin{enumerate}
  \item \kw{商品与服务}的跨境交易（transactions of goods and services across borders）；
  \item \kw{要素（资本、劳动等）}的跨境流动（flows of factors）；
  \item \kw{生产过程}的跨境重新配置（reallocation of production process，如外包、全球价值链）；
  \item 针对上述活动的\term{政策}（关税、补贴、贸易协定等）。
\end{enumerate}

\begin{intubox}[课堂手写批注（第 7 页）]
老师圈出 ``International'' 一词并手写批注：``$\Rightarrow$ extension, regional, spatial''——
即``国际''的实质是\textbf{跨越边界的空间延伸}，``国际贸易''中的``Trade''指商品/要素/生产的交易，``International''强调其跨越边界的空间属性（regional、spatial extension）。
\end{intubox}

\textbf{国际贸易研究的三大类问题（第 8--9 页）}：
\begin{enumerate}
  \item \term{为什么各国会贸易？（Why do countries trade?）}
    \begin{itemize}
      \item \kw{禀赋（endowments）}差异——商品与要素的数量差异；
      \item \term{技术}差异；
      \item \term{规模经济}（economies of scale）。
    \end{itemize}
    这三个原因分别对应本课程模块一中的各个模型（禀赋模型、李嘉图模型、H-O 模型、规模经济与新贸易理论）。
  \item \term{贸易的后果是什么？（What are the consequences of trade?）}
    \begin{itemize}
      \item 工资、价格、就业与\kw{福利}（老师手写补充：``(price, quantities)''——福利分析最终落到价格与数量上）；
      \item 可获得的商品的\kw{质量与种类}（quality and variety）；
      \item 投入品（inputs）的质量与种类；
      \item \kw{赢家与输家（winners and losers）}：工人、企业——贸易有总体收益但存在分配效应。
    \end{itemize}
  \item \term{影响贸易的政策（Policies that affect trade）}：
    \begin{itemize}
      \item 工具：关税（tariffs）、反倾销税（antidumping duties）、出口补贴（export subsidies）、反补贴税（countervailing duties）、监管壁垒（regulatory barriers）；
      \item \kw{贸易壁垒的成因与后果}（老师红笔下划线强调）；
      \item 多边贸易协定——WTO；区域/优惠贸易协定——EU、NAFTA、中国--东盟。
    \end{itemize}
\end{enumerate}

\secn{1.2 贸易研究的方法论（全课程总纲）}

\begin{mapbox}[模型构建方法地图（通用框架，第 10--16 页）]
\textbf{本讲的方法论地图}——这是课程中每一个贸易模型共享的``骨架''，复习每个模型时都应回到这张图：
\begin{enumerate}
  \item \term{模型（Models）}：用一系列模型以严格的方式传达非平凡、常常出人意料的结论；有时会超越传统智慧与直觉。
  \item \term{模型的结构（structure）}：
    \[
    \text{\kw{外生变量}} \;\longrightarrow\; \text{\term{模型}} \;\longrightarrow\; \text{\kw{内生变量}}
    \]
    外生变量（exogenous）来自模型之外（给定的禀赋、技术、偏好、政策等）；内生变量（endogenous）是模型要解释的对象（价格、产量、贸易模式、福利等）。模型说明外生变量的变化如何影响内生变量（见图~\ref{fig:ch01-modelwork}）。
  \item \kw{比较静态（Comparative Statics）}：改变\textbf{一个}外生变量、保持其他不变，内生变量如何变化？——这是每讲分析的核心动作。
  \item \term{假设（Assumptions）}：``All models are wrong, but some are useful!''（所有模型都是错的，但有些是有用的。）对每个模型要问：\textbf{关键假设是什么？什么时候该假设成立/不成立？若打破该假设，结论是否还成立？}
  \item \term{两类均衡方法}：
    \begin{itemize}
      \item \term{局部均衡（partial equilibrium）}：一次只分析\textbf{一个市场}的供求与均衡；
      \item \kw{一般均衡（general equilibrium）}：同时分析\textbf{多个市场}的供求与均衡。
    \end{itemize}
    本课程\textbf{以一般均衡为主}——因为贸易问题通常需要同时分析不同产品的市场、不同国家的市场。
  \item \term{经验研究（Empirics）}：模型给出重要预测后，用现实数据检验模型的关键关系；评估现实冲击与政策。本课程约 75\% 时间讲模型、25\% 时间讲经验证据。
\end{enumerate}
\end{mapbox}

\begin{figure}[H]
  \centering
  \includegraphics[width=0.72\linewidth]{figs/ch01_fig_1.png}
  \caption{模型如何工作（Mankiw, \textit{Macroeconomics} Figure 1-4）：外生变量 $\to$ 模型 $\to$ 内生变量。}
  \label{fig:ch01-modelwork}
\end{figure}

\begin{exbox}[一般均衡思维的热身：``破窗谬论''（broken window fallacy, Bastiat）]
杂货店的玻璃窗被顽童打破。``聪明人''说：这\textbf{对经济是好事}——店主必须付钱请镇上的玻璃匠换窗，玻璃匠又会把这笔额外的钱花出去，从而拉动本地经济。

\textbf{为什么``聪明人''不完全正确}（PPT 第 15 页蓝字）：
\begin{itemize}
  \item 影响力传导链（``chain of influence''）中存在\term{摩擦}（传导不完全）；
  \item ``负向影响链''很可能相互抵消，最终剩下的净效应只是\textbf{碎了一块玻璃}——资源本可用于其他用途（机会成本）。
\end{itemize}
\textbf{与贸易的联系}：只看到可见渠道（玻璃匠得益）而忽略全部市场的一般均衡反馈，正是破窗谬论的问题；评估贸易政策同样需要一般均衡思维。
\end{exbox}

\begin{intubox}[先修要求（第 22--23 页，复习自检清单）]
\textbf{微观经济学}：PPF、机会成本、供求曲线及其移动、比较静态；消费者理论（偏好、效用、无差异曲线、预算约束、最优化）；生产者理论（生产函数、MPL/MPK、规模报酬、要素供求与要素价格、完全竞争利润最大化、垄断竞争）；福利分析（消费者剩余、生产者剩余、无谓损失）；\kw{一般均衡}：Edgeworth 盒、瓦尔拉斯定理（Walras' Law）。

\textbf{手写批注}（第 22 页，老师手绘 Edgeworth 盒草图见图~\ref{fig:ch01-edgeworth}）：
``$N$ 个市场：若第 $1,\dots,N-1$ 个市场都均衡 $\Rightarrow$ 第 $N$ 个市场自动均衡''——即\term{瓦尔拉斯定律}：只要前 $N-1$ 个市场出清，价值加总意味着第 $N$ 个市场必然出清，自由度只剩 $N-1$ 个。

\textbf{数学}：基本代数；（偏）导数与全微分；积分的概念；比例变化（``弹性''）：
\[{\color{hlorange}
y = x_1^{\alpha_1} x_2^{\alpha_2} \cdots x_N^{\alpha_N}}
\;\Longrightarrow\;
\frac{\Delta y}{y} \approx \alpha_1\frac{\Delta x_1}{x_1} + \alpha_2\frac{\Delta x_2}{x_2} + \cdots + \alpha_N\frac{\Delta x_N}{x_N}
\]
（Cobb--Douglas 型函数的比例变化法则，PPT 荧光高亮；对数线性化后即弹性加权和）；拉格朗日式优化：$\max_{C_1,C_2} U(C_1,C_2) \;\text{s.t.}\; p_1 C_1 + p_2 C_2 \le I$。
\end{intubox}

\begin{figure}[H]
  \centering
  \includegraphics[width=0.3\linewidth]{figs/ch01_fig_8.png}
  \caption{老师板书：Edgeworth 盒草图——纵轴为资本 $K$、横轴为劳动 $L$，两个国家 $A$、$B$ 各自占有要素组合 $(K_A, L_A)$ 与 $(K_B, L_B)$（用于讲解一般均衡与瓦尔拉斯定律）。}
  \label{fig:ch01-edgeworth}
\end{figure}

\secn{1.3 课程结构与考核（速览）}

\textbf{课程四大模块（第 17--20 页）}：
\begin{itemize}
  \item \term{模块一：古典贸易理论}——比较优势及其决定因素、贸易的福利影响；\term{Endowment Model（禀赋模型）、Ricardian Model（李嘉图模型）、Specific Factors Model（特定要素模型）、Heckscher--Ohlin Model（H-O 模型）}；
  \item \term{模块二：新贸易理论的``玩具模型''（\kw{new trade theory}，荧光高亮）}——要素与生产过程的流动、外部与内部规模经济下的贸易、不完全竞争下的贸易、异质性企业下的贸易；
  \item \term{模块三：贸易政策}——贸易政策如何影响贸易模式？贸易政策的福利影响？什么政治经济因素决定贸易政策？
  \item \term{模块四*：贸易的经验研究}——经验研究导论、比较优势的经验证据、贸易如何影响微观主体、贸易成本与经济地理、中美贸易问题、贸易与环境。
\end{itemize}

\textbf{考核（第 24--27 页）}：出勤 10\%（扫码，可缺勤两次）；课堂参与 10\%（经验文献课前读摘要与引言，课堂小测）；习题集 15\%（4 次，check-plus/check/check-minus，可弃一次，迟交不收）；期中 25\%（11 月 6 日，课内）；期末 30\%（1 月 6 日，覆盖全学期，约 2/3 内容在期中之后）；文献综述 10\%（12 月 6 日前，2 页）；加分作业（额外 2\%）：研究计划（3 页内，1--3 人组队），截止日期\textbf{12 月 6 日}（老师手写将打印版的 Dec.~3 改为 Dec.~6）。所有考试开卷、可带计算器。

\textbf{参考教材}（非强制，作业与考试以课堂讲义为准）：Krugman, Obstfeld \& Melitz, \textit{International Economics: Theory and Policy}, 11th ed.；Feenstra \& Taylor, \textit{International Trade}, 3rd ed.

\secn{1.4 国际贸易的五个流行误解与事实}

\begin{center}
\begin{tabular}{@{}p{0.30\linewidth}p{0.66\linewidth}@{}}
\toprule
\textbf{误解（Misconception）} & \textbf{事实（Reality）} \\
\midrule
1. ``我们生活在一个国际贸易占主导地位的全球化世界'' &
贸易相对于国内交易并没有那么多：多数国家的\term{贸易/GDP 比率}很小（见表~\ref{tab:ch01-ratio}）；但过去几十年贸易量持续上升——``全球化''。 \\
2. ``我们主要贸易消费品，如食品或 iPhone'' &
\term{大部分贸易是工业用品与资本品}（industrial supplies and capital goods，如机械、中间投入品）。 \\
3. ``贸易（和 FDI）主要发生在高工资（富国）与低工资（穷国）之间'' &
\term{大部分贸易发生在富国之间}；FDI 的来源与去向也都是富国为主。 \\
4. ``世界已经变平，距离不再影响贸易模式'' &
\term{贸易对距离反而更敏感了}：距离对贸易的负效应在过去几十年间增大（见图~\ref{fig:ch01-distance}）。 \\
5. ``贸易是零和博弈，有赢家必有输家'' &
\term{贸易是正和博弈}：所有国家都从贸易中获益（国家层面）；但\textbf{一国内部}既有赢家也有输家。这可能是关于贸易最重要的误解——后续课程将用模型与数据展开（追问：既然如此，为什么仍有保护主义？）。 \\
\bottomrule
\end{tabular}
\end{center}

\begin{table}[H]
  \centering
  \caption{2014 年部分经济体贸易/GDP 比率（(进口+出口)/GDP，\%）；数据来源：World Development Indicators, the World Bank（PPT 第 36 页）。}
  \label{tab:ch01-ratio}
  \begin{tabular}{@{}lrr@{}}
    \toprule
    经济体 & (进口+出口)/GDP & GDP（十亿美元） \\
    \midrule
    中国香港 & 220 & 291 \\
    新加坡 & 117 & 308 \\
    德国 & 42 & 3{,}853 \\
    俄罗斯 & 41 & 1{,}861 \\
    加拿大 & 32 & 1{,}787 \\
    英国 & 29 & 2{,}942 \\
    印度 & 25 & 2{,}067 \\
    中国 & 21 & 10{,}360 \\
    日本 & 18 & 4{,}920 \\
    美国 & 15 & 16{,}768 \\
    \bottomrule
  \end{tabular}
\end{table}

\begin{intubox}[全球化的三个``黄金时代''与动因（第 37--41 页）]
贸易/GDP 比率的历史（图~\ref{fig:ch01-history}）显示三个``黄金时代''：\kw{1890s 至一战}、\kw{1950 年代以来}、\kw{1990 年代以来}——全球化并非当代独有，历史上已有过一波（19 世纪末的全球化在两次大战间倒退）。

\textbf{塑造全球化的力量}：(1) \term{技术}——交通改善（transportation-enhancing）、通讯改善（communication-enhancing）、生产率提升（productivity-enhancing）；(2) \term{政策}——贸易自由化、FDI 自由化、国际移民，或者相反（转向保护主义）。

\textbf{最近这波全球化的新特点}：贸易与资本流动规模空前；通讯技术；\term{生产的碎片化}（fragmentation of production，如 iPhone 全球生产链）；保护主义回潮（protectionism bouncing back）。
\end{intubox}

\begin{figure}[H]
  \centering
  \includegraphics[width=0.6\linewidth]{figs/ch01_fig_2.png}
  \caption{贸易/GDP 比率的历史（1890--2011，澳大利亚、加拿大、欧洲、日本、英国、美国）：三个``黄金时代''清晰可见（PPT 第 38 页）。}
  \label{fig:ch01-history}
\end{figure}

\begin{figure}[H]
  \centering
  \includegraphics[width=0.88\linewidth]{figs/ch01_fig_3.png}
  \caption{美国进口结构（1925--2014）：工业用品与材料占比大幅下降，资本品与消费品上升——大部分贸易并非消费品。来源：Bureau of Economic Analysis（PPT 第 43 页）。}
  \label{fig:ch01-imports}
\end{figure}

\begin{figure}[H]
  \centering
  \includegraphics[width=0.88\linewidth]{figs/ch01_fig_4.png}
  \caption{美国出口结构（1925--2014）：资本品约占 60\%、工业用品约占 40\%，食品饲料饮料占比很小（PPT 第 44 页）。}
  \label{fig:ch01-exports}
\end{figure}

\begin{figure}[H]
  \centering
  \includegraphics[width=0.92\linewidth]{figs/ch01_fig_5.png}
  \caption{2014 年世界贸易地图（商品贸易流，总额 \$19{,}790 十亿）：欧洲内部贸易 \$4{,}566 十亿为全球最大，其次是亚洲内部 \$2{,}785 十亿；最大单一流向为亚洲$\to$欧洲 \$1{,}032 十亿。来源：United Nations trade data（PPT 第 46 页）。}
  \label{fig:ch01-trademap}
\end{figure}

\begin{figure}[H]
  \centering
  \includegraphics[width=0.75\linewidth]{figs/ch01_fig_6.png}
  \caption{2013 年世界移民流地图（总量 2.32 亿人）：与贸易地图同样呈现``富国之间为主''的格局（PPT 第 47 页）。}
  \label{fig:ch01-fdimap}
\end{figure}

\begin{figure}[H]
  \centering
  \includegraphics[width=0.8\linewidth]{figs/ch01_fig_7.png}
  \caption{``距离之谜''：引力方程中距离系数的 OLS 估计随年份变化（Yotov 2012, \textit{Economics Letters} 117）——估计值从 1965 年的约 $-1$ 增大到 2005 年的约 $-1.7$，贸易对距离越来越敏感（PPT 第 49 页）。}
  \label{fig:ch01-distance}
\end{figure}

\begin{intubox}[课堂手写批注：引力方程（第 49 页）]
老师在标题下手写``Gravity''：$F \propto m_1 m_2 r^{-2}$（牛顿万有引力），并类推贸易的引力方程
\[
X_{ij} \propto Y_i\, Y_j\, \mathrm{Dist}_{ij}^{-1}
\]
（$X_{ij}$ 为 $i$、$j$ 两国间贸易额，$Y_i, Y_j$ 为两国经济规模，$\mathrm{Dist}_{ij}$ 为距离）；图中纵轴的距离估计系数旁手写 $\hat\lambda$，即上式中的距离弹性（为负且绝对值在增大）。
\end{intubox}

\begin{intubox}[距离代理了什么？（第 50 页）]
距离不只是物理里程，它代理了四大类``贸易成本''：
\begin{itemize}
  \item \term{文化（Cultural）}：语言、族群、宗教、规范等；
  \item \term{行政（Administrative）}：殖民纽带、监管、敌对关系、汇率制度等；
  \item \term{地理（Geographic）}：物理距离、共同边界、港口可及性、气候等；
  \item \term{经济（Economic）}：禀赋、市场规模、汇率波动、基础设施等。
\end{itemize}
\end{intubox}

\begin{sumbox}[本讲总结：通用概念与方法]
\begin{enumerate}
  \item \term{核心定义}：出口/进口、贸易顺差/逆差（净出口的市场价值）、FDI（对海外关联企业的投资与控制）；本课程不含汇率与国际金融。
  \item \term{贸易的三个层次}：商品与服务交易 $\to$ 要素流动（资本、劳动）$\to$ 生产过程跨境重新配置（外包/全球价值链）；外加贸易政策。
  \item \term{贸易的原因}：禀赋差异、技术差异、规模经济——分别对应禀赋模型、李嘉图模型、新贸易理论。
  \item \term{贸易的后果}：价格与数量（进而福利）、产品与投入品的种类与质量、国家内部的赢家与输家（分配效应）。
  \item \term{模型方法论}（每个模型的骨架）：外生变量 $\to$ 模型 $\to$ 内生变量；比较静态（只动一个外生变量）；重视关键假设及其稳健性（``All models are wrong, but some are useful''）；以\textbf{一般均衡}为主要分析框架（多市场同时均衡）；模型预测需要经验检验。
  \item \term{五大事实}：贸易/GDP 不高但在增长（三个黄金时代）；贸易品以工业用品与资本品为主（中间品贸易为重要现象）；贸易与 FDI 主要发生在富国之间；距离的负效应在增强（引力方程 $X_{ij} \propto Y_i Y_j \mathrm{Dist}_{ij}^{-1}$）；贸易是国家层面正和、国内层面存在分配冲突的博弈。
  \item \term{瓦尔拉斯定律}（板书）：$N$ 个市场中前 $N-1$ 个均衡，则第 $N$ 个自动均衡——一般均衡分析中独立的市场约束只有 $N-1$ 个。
\end{enumerate}
\end{sumbox}
